% chapterintro
% ==============================
% 第2章 导言：2.0-chinese.md
本章的主要内容有：

\begin{itemize}
\item 检测操作系统

\item 处理与平台相关的源码

\item 处理与编译器相关的源码

\item 检测处理器体系结构

\item 检测处理器指令集

\item 为Eigen库使能向量化

\end{itemize}

尽管CMake跨平台，但有时源代码并不是完全可移植(例如：当使用依赖于供应商的扩展时)，我们努力使源代码能够跨平台、操作系统和编译器。这个过程中会发现，有必要根据平台不同的方式配置和/或构建代码。这对于历史代码或交叉编译尤其重要，我们将在第13章中讨论这个主题。了解处理器指令集也有助于优化特定目标平台的性能。本章会介绍，检测环境的方法，并给出建议。


% --- recipe 1: 2.1-chinese.md ---
\section{2.1 检测操作系统}

\begin{myInfoBox}
\textbf{NOTE}：此示例代码可以在 \texthref{https://github.com/dev-cafe/cmake-cookbook/tree/v1.0/chapter-02/recipe-01} 中找到。该示例在CMake 3.5版(或更高版本)中是有效的，并且已经在GNU/Linux、macOS和Windows上进行过测试。
\end{myInfoBox}

CMake是一组跨平台工具。不过，了解操作系统(OS)上执行配置或构建步骤也很重要。从而与操作系统相关的CMake代码，会根据操作系统启用条件编译，或者在可用或必要时使用特定于编译器的扩展。本示例中，我们将通过一个不需要编译任何源代码的示例，演示如何使用CMake检测操作系统。为了简单起见，我们只考虑配置过程。

\mySmallsection{具体实施}

我们将用一个非常简单的\texttt{CMakeLists.txt}进行演示:

\begin{enumerate}
\item 首先，定义CMake最低版本和项目名称。请注意，语言是\texttt{NONE}:

\begin{codebox}[cmake]
cmake_minimum_required(VERSION 3.5 FATAL_ERROR)
project(recipe-01 LANGUAGES NONE)
\end{codebox}

\item 然后，根据检测到的操作系统信息打印消息:

\begin{codebox}[cmake]
if(CMAKE_SYSTEM_NAME STREQUAL "Linux")
    message(STATUS "Configuring on/for Linux")
elseif(CMAKE_SYSTEM_NAME STREQUAL "Darwin")
    message(STATUS "Configuring on/for macOS")
elseif(CMAKE_SYSTEM_NAME STREQUAL "Windows")
    message(STATUS "Configuring on/for Windows")
elseif(CMAKE_SYSTEM_NAME STREQUAL "AIX")
    message(STATUS "Configuring on/for IBM AIX")
else()
    message(STATUS "Configuring on/for ${CMAKE_SYSTEM_NAME}")
endif()
\end{codebox}

测试之前，检查前面的代码块，并考虑相应系统上的具体行为。

\item 现在，测试配置项目:

\begin{codebox}[shell]
$ mkdir -p build
$ cd build
$ cmake ..
\end{codebox}

\item 关于CMake输出，这里有一行很有趣——在Linux系统上(在其他系统上，输出会不同):

\begin{codebox}[shell]
-- Configuring on/for Linux
\end{codebox}

\end{enumerate}

\mySmallsection{工作原理}

CMake为目标操作系统定义了\texttt{CMAKE\_SYSTEM\_NAME}，因此不需要使用定制命令、工具或脚本来查询此信息。然后，可以使用此变量的值实现特定于操作系统的条件和解决方案。在具有\texttt{uname}命令的系统上，将此变量设置为\texttt{uname -s}的输出。该变量在macOS上设置为“Darwin”。在Linux和Windows上，它分别计算为“Linux”和“Windows”。我们了解了如何在特定的操作系统上执行特定的CMake代码。当然，应该尽量减少这种定制化行为，以便简化迁移到新平台的过程。

\begin{myInfoBox}
\textbf{NOTE}：为了最小化从一个平台转移到另一个平台时的成本，应该避免直接使用Shell命令，还应该避免显式的路径分隔符(Linux和macOS上的前斜杠和Windows上的后斜杠)。CMake代码中只使用前斜杠作为路径分隔符，CMake将自动将它们转换为所涉及的操作系统环境。
\end{myInfoBox}


% --- recipe 2: 2.2-chinese.md ---
\section{2.2 处理与平台相关的源代码}

\begin{myInfoBox}
\textbf{NOTE}：此示例代码可以在 \texthref{https://github.com/dev-cafe/cmake-cookbook/tree/v1.0/chapter-02/recipe-02} 中找到，包含一个C++示例。该示例在CMake 3.5版(或更高版本)中是有效的，并且已经在GNU/Linux、macOS和Windows上进行过测试。
\end{myInfoBox}

理想情况下，应该避免依赖于平台的源代码，但是有时我们没有选择，特别是当要求配置和编译不是自己编写的代码时。本示例中，将演示如何使用CMake根据操作系统编译源代码。

\mySmallsection{准备工作}

修改\texttt{hello-world.cpp}示例代码，将第1章第1节的例子进行修改:

\begin{codebox}[cpp]
#include <cstdlib>
#include <iostream>
#include <string>

std::string say_hello() {
#ifdef IS_WINDOWS
  return std::string("Hello from Windows!");
#elif IS_LINUX
  return std::string("Hello from Linux!");
#elif IS_MACOS
  return std::string("Hello from macOS!");
#else
  return std::string("Hello from an unknown system!");
#endif
}

int main() {
  std::cout << say_hello() << std::endl;
  return EXIT_SUCCESS;
}
\end{codebox}

\mySmallsection{具体实施}

完成一个\texttt{CMakeLists.txt}实例，使我们能够基于目标操作系统有条件地编译源代码：

\begin{enumerate}
\item 首先，设置了CMake最低版本、项目名称和支持的语言:

\begin{codebox}[cmake]
cmake_minimum_required(VERSION 3.5 FATAL_ERROR)
project(recipe-02 LANGUAGES CXX)
\end{codebox}

\item 然后，定义可执行文件及其对应的源文件:

\begin{codebox}[cmake]
add_executable(hello-world hello-world.cpp)
\end{codebox}

\item 通过定义以下目标编译定义，让预处理器知道系统名称:

\begin{codebox}[cmake]
if(CMAKE_SYSTEM_NAME STREQUAL "Linux")
  target_compile_definitions(hello-world PUBLIC "IS_LINUX")
endif()
if(CMAKE_SYSTEM_NAME STREQUAL "Darwin")
  target_compile_definitions(hello-world PUBLIC "IS_MACOS")
endif()
if(CMAKE_SYSTEM_NAME STREQUAL "Windows")
  target_compile_definitions(hello-world PUBLIC "IS_WINDOWS")
endif()
\end{codebox}

继续之前，先检查前面的表达式，并考虑在不同系统上有哪些行为。

\item 现在，准备测试它，并配置项目:

\begin{codebox}[shell]
$ mkdir -p build
$ cd build
$ cmake ..
$ cmake --build .
$ ./hello-world
Hello from Linux!
\end{codebox}

\end{enumerate}

Windows系统上，将看到来自Windows的Hello。其他操作系统将产生不同的输出。

\mySmallsection{工作原理}

\texttt{hello-world.cpp}示例中，有趣的部分是基于预处理器定义\texttt{IS\_WINDOWS}、\texttt{IS\_LINUX}或\texttt{IS\_MACOS}的条件编译:

\begin{codebox}[cmake]
std::string say_hello() {
#ifdef IS_WINDOWS
  return std::string("Hello from Windows!");
#elif IS_LINUX
  return std::string("Hello from Linux!");
#elif IS_MACOS
  return std::string("Hello from macOS!");
#else
  return std::string("Hello from an unknown system!");
#endif
}
\end{codebox}

这些定义在CMakeLists.txt中配置时定义，通过使用\texttt{target\_compile\_definition}在预处理阶段使用。可以不重复\texttt{if-endif}语句，以更紧凑的表达式实现，我们将在下一个示例中演示这种重构方式。也可以把\texttt{if-endif}语句加入到一个\texttt{if-else-else-endif}语句中。这个阶段，可以使用\texttt{add\_definitions(-DIS\_LINUX)}来设置定义(当然，可以根据平台调整定义)，而不是使用\texttt{target\_compile\_definition}。使用\texttt{add\_definitions}的缺点是，会修改编译整个项目的定义，而\texttt{target\_compile\_definitions}给我们机会，将定义限制于一个特定的目标，以及通过\texttt{ PRIVATE|PUBLIC|INTERFACE }限定符，限制这些定义可见性。第1章的第8节，对这些限定符有详细的说明:

\begin{itemize}
\item \textbf{PRIVATE}，编译定义将只应用于给定的目标，而不应用于相关的其他目标。

\item \textbf{INTERFACE}，对给定目标的编译定义将只应用于使用它的目标。

\item \textbf{PUBLIC}，编译定义将应用于给定的目标和使用它的所有其他目标。

\end{itemize}

\begin{myInfoBox}
\textbf{NOTE}：将项目中的源代码与平台相关性最小化，可使移植更加容易。
\end{myInfoBox}


% --- recipe 3: 2.3-chinese.md ---
\section{2.3 处理与编译器相关的源代码}

\begin{myInfoBox}
\textbf{NOTE}：此示例代码可以在 \texthref{https://github.com/dev-cafe/cmake-cookbook/tree/v1.0/chapter-02/recipe-03} 中找到，包含一个C++和Fortran示例。该示例在CMake 3.5版(或更高版本)中是有效的，并且已经在GNU/Linux、macOS和Windows上进行过测试。
\end{myInfoBox}

这个方法与前面的方法类似，我们将使用CMake来编译依赖于环境的条件源代码：本例将依赖于编译器。为了可移植性，我们尽量避免去编写新代码，但遇到有依赖的情况我们也要去解决，特别是当使用历史代码或处理编译器依赖工具，如\href{https://github.com/google/sanitizers}{sanitizers}。从这一章和前一章的示例中，我们已经掌握了实现这一目标的所有方法。尽管如此，讨论与编译器相关的源代码的处理问题还是很有用的，这样我们将有机会从另一方面了解CMake。

\mySmallsection{准备工作}

本示例中，我们将从\texttt{C++}中的一个示例开始，稍后我们将演示一个\texttt{Fortran}示例，并尝试重构和简化CMake代码。

看一下\texttt{hello-world.cpp}源代码:

\begin{codebox}[cpp]
#include <cstdlib>
#include <iostream>
#include <string>

std::string say_hello() {
#ifdef IS_INTEL_CXX_COMPILER
  // only compiled when Intel compiler is selected
  // such compiler will not compile the other branches
  return std::string("Hello Intel compiler!");
#elif IS_GNU_CXX_COMPILER
  // only compiled when GNU compiler is selected
  // such compiler will not compile the other branches
  return std::string("Hello GNU compiler!");
#elif IS_PGI_CXX_COMPILER
  // etc.
  return std::string("Hello PGI compiler!");
#elif IS_XL_CXX_COMPILER
  return std::string("Hello XL compiler!");
#else
  return std::string("Hello unknown compiler - have we met before?");
#endif
}

int main() {
  std::cout << say_hello() << std::endl;
  std::cout << "compiler name is " COMPILER_NAME << std::endl;
  return EXIT_SUCCESS;
}
\end{codebox}

\texttt{Fortran}示例(\texttt{hello-world.F90}):

\begin{codebox}[fortran]
program hello

  implicit none
#ifdef IS_Intel_FORTRAN_COMPILER
  print *, 'Hello Intel compiler!'
#elif IS_GNU_FORTRAN_COMPILER
  print *, 'Hello GNU compiler!'
#elif IS_PGI_FORTRAN_COMPILER
  print *, 'Hello PGI compiler!'
#elif IS_XL_FORTRAN_COMPILER
  print *, 'Hello XL compiler!'
#else
  print *, 'Hello unknown compiler - have we met before?'
#endif

end program
\end{codebox}

\mySmallsection{具体实施}

我们将从\texttt{C++}的例子开始，然后再看\texttt{Fortran}的例子:

\begin{enumerate}
\item \texttt{CMakeLists.txt}文件中，定义了CMake最低版本、项目名称和支持的语言:

\begin{codebox}[cmake]
cmake_minimum_required(VERSION 3.5 FATAL_ERROR)
project(recipe-03 LANGUAGES CXX)
\end{codebox}

\item 然后，定义可执行目标及其对应的源文件:

\begin{codebox}[cmake]
add_executable(hello-world hello-world.cpp)
\end{codebox}

\item 通过定义以下目标编译定义，让预处理器了解编译器的名称和供应商:

\begin{codebox}[cmake]
target_compile_definitions(hello-world PUBLIC "COMPILER_NAME=\"${CMAKE_CXX_COMPILER_ID}\"")
if(CMAKE_CXX_COMPILER_ID MATCHES Intel)
  target_compile_definitions(hello-world PUBLIC "IS_INTEL_CXX_COMPILER")
endif()
if(CMAKE_CXX_COMPILER_ID MATCHES GNU)
  target_compile_definitions(hello-world PUBLIC "IS_GNU_CXX_COMPILER")
endif()
if(CMAKE_CXX_COMPILER_ID MATCHES PGI)
  target_compile_definitions(hello-world PUBLIC "IS_PGI_CXX_COMPILER")
endif()
if(CMAKE_CXX_COMPILER_ID MATCHES XL)
  target_compile_definitions(hello-world PUBLIC "IS_XL_CXX_COMPILER")
endif()
\end{codebox}

\end{enumerate}

现在我们已经可以预测结果了:

\begin{codebox}[shell]
$ mkdir -p build
$ cd build
$ cmake ..
$ cmake --build .
$ ./hello-world

Hello GNU compiler!
\end{codebox}

使用不同的编译器，此示例代码将打印不同的问候语。

前一个示例的\texttt{CMakeLists.txt}文件中的\texttt{if}语句似乎是重复的，我们不喜欢重复的语句。能更简洁地表达吗？当然可以！为此，让我们再来看看\texttt{Fortran}示例。

\texttt{Fortran}例子的\texttt{CMakeLists.txt}文件中，我们需要做以下工作:

\begin{enumerate}
\item 需要使\texttt{Fortran}语言:

\begin{codebox}[cmake]
project(recipe-03 LANGUAGES Fortran)
\end{codebox}

\item 然后，定义可执行文件及其对应的源文件。在本例中，使用大写\texttt{.F90}后缀:

\begin{codebox}[cmake]
add_executable(hello-world hello-world.F90)
\end{codebox}

\item 我们通过定义下面的目标编译定义，让预处理器非常清楚地了解编译器:

\begin{codebox}[cmake]
target_compile_definitions(hello-world
  PUBLIC "IS_${CMAKE_Fortran_COMPILER_ID}_FORTRAN_COMPILER"
  )
\end{codebox}

\end{enumerate}

其余行为与\texttt{C++}示例相同。

\mySmallsection{工作原理}

\texttt{CMakeLists.txt}会在配置时，进行预处理定义，并传递给预处理器。\texttt{Fortran}示例包含非常紧凑的表达式，我们使用\texttt{CMAKE\_Fortran\_COMPILER\_ID}变量，通过\texttt{target\_compile\_definition}使用构造预处理器进行预处理定义。为了适应这种情况，我们必须将"Intel"从\texttt{IS\_INTEL\_CXX\_COMPILER}更改为\texttt{IS\_Intel\_FORTRAN\_COMPILER}。通过使用相应的\texttt{CMAKE\_C\_COMPILER\_ID}和\texttt{CMAKE\_CXX\_COMPILER\_ID}变量，我们可以在\texttt{C}或\texttt{C++}中实现相同的效果。但是，请注意，\texttt{Fortran}0不能保证为所有编译器或语言都定义。

\begin{myInfoBox}
\textbf{NOTE}：对于应该预处理的\texttt{Fortran}代码使用\texttt{.F90}后缀，对于不需要预处理的代码使用\texttt{.f90}后缀。
\end{myInfoBox}


% --- recipe 4: 2.4-chinese.md ---
\section{2.4 检测处理器体系结构}

\begin{myInfoBox}
\textbf{NOTE}：此示例代码可以在 \texthref{https://github.com/dev-cafe/cmake-cookbook/tree/v1.0/chapter-02/recipe-04} 中找到，包含一个C++示例。该示例在CMake 3.5版(或更高版本)中是有效的，并且已经在GNU/Linux、macOS和Windows上进行过测试。
\end{myInfoBox}

19世纪70年代，出现的64位整数运算和本世纪初出现的用于个人计算机的64位寻址，扩大了内存寻址范围，开发商投入了大量资源来移植为32位体系结构硬编码，以支持64位寻址。许多博客文章，如 \texthref{https://www.viva64.com/en/a/0004/} ，致力于讨论将\texttt{C++}代码移植到64位平台中的典型问题和解决方案。虽然，避免显式硬编码的方式非常明智，但需要在使用CMake配置的代码中适应硬编码限制。本示例中，我们会来讨论检测主机处理器体系结构的选项。

\mySmallsection{准备工作}

我们以下面的\texttt{arch-dependent.cpp}代码为例：

\begin{codebox}[cpp]
#include <cstdlib>
#include <iostream>
#include <string>

#define STRINGIFY(x) #x
#define TOSTRING(x) STRINGIFY(x)

std::string say_hello()
{
  std::string arch_info(TOSTRING(ARCHITECTURE));
  arch_info += std::string(" architecture. ");
#ifdef IS_32_BIT_ARCH
  return arch_info + std::string("Compiled on a 32 bit host processor.");
#elif IS_64_BIT_ARCH
  return arch_info + std::string("Compiled on a 64 bit host processor.");
#else
  return arch_info + std::string("Neither 32 nor 64 bit, puzzling ...");
#endif
}

int main()
{
  std::cout << say_hello() << std::endl;
  return EXIT_SUCCESS;
}
\end{codebox}

\mySmallsection{具体实施}

\texttt{CMakeLists.txt}文件中，我们需要以下内容:

\begin{enumerate}
\item 首先，定义可执行文件及其源文件依赖关系:

\begin{codebox}[cmake]
cmake_minimum_required(VERSION 3.5 FATAL_ERROR)
project(recipe-04 LANGUAGES CXX)
add_executable(arch-dependent arch-dependent.cpp)
\end{codebox}

\item 检查空指针类型的大小。CMake的\texttt{CMAKE\_SIZEOF\_VOID\_P}变量会告诉我们CPU是32位还是64位。我们通过状态消息让用户知道检测到的大小，并设置预处理器定义:

\begin{codebox}[cmake]
if(CMAKE_SIZEOF_VOID_P EQUAL 8)
  target_compile_definitions(arch-dependent PUBLIC "IS_64_BIT_ARCH")
  message(STATUS "Target is 64 bits")
else()
  target_compile_definitions(arch-dependent PUBLIC "IS_32_BIT_ARCH")
  message(STATUS "Target is 32 bits")
endif()
\end{codebox}

\item 通过定义以下目标编译定义，让预处理器了解主机处理器架构，同时在配置过程中打印状态消息:

\begin{codebox}[cmake]
if(CMAKE_HOST_SYSTEM_PROCESSOR MATCHES "i386")
    message(STATUS "i386 architecture detected")
elseif(CMAKE_HOST_SYSTEM_PROCESSOR MATCHES "i686")
    message(STATUS "i686 architecture detected")
elseif(CMAKE_HOST_SYSTEM_PROCESSOR MATCHES "x86_64")
    message(STATUS "x86_64 architecture detected")
else()
    message(STATUS "host processor architecture is unknown")
endif()
target_compile_definitions(arch-dependent
  PUBLIC "ARCHITECTURE=${CMAKE_HOST_SYSTEM_PROCESSOR}"
  )
\end{codebox}

\item 配置项目，并注意状态消息(打印出的信息可能会发生变化):

\begin{codebox}[shell]
$ mkdir -p build
$ cd build
$ cmake ..
...
-- Target is 64 bits
-- x86_64 architecture detected
...
\end{codebox}

\item 最后，构建并执行代码(实际输出将取决于处理器架构):

\begin{codebox}[shell]
$ cmake --build .
$ ./arch-dependent
x86_64 architecture. Compiled on a 64 bit host processor.
\end{codebox}

\end{enumerate}

\mySmallsection{工作原理}

CMake定义了\texttt{CMAKE\_HOST\_SYSTEM\_PROCESSOR}变量，以包含当前运行的处理器的名称。可以设置为“i386”、“i686”、“x86\_64”、“AMD64”等等，当然，这取决于当前的CPU。\texttt{CMAKE\_SIZEOF\_VOID\_P}为void指针的大小。我们可以在CMake配置时进行查询，以便修改目标或目标编译定义。可以基于检测到的主机处理器体系结构，使用预处理器定义，确定需要编译的分支源代码。正如在前面的示例中所讨论的，编写新代码时应该避免这种依赖，但在处理遗留代码或交叉编译时，这种依赖是有用的，交叉编译会在第13章进行讨论。

\begin{myInfoBox}
\textbf{NOTE}：使用\texttt{CMAKE\_SIZEOF\_VOID\_P}是检查当前CPU是否具有32位或64位架构的唯一“真正”可移植的方法。
\end{myInfoBox}

\mySmallsection{更多信息}

除了\texttt{CMAKE\_HOST\_SYSTEM\_PROCESSOR}, CMake还定义了\texttt{CMAKE\_SYSTEM\_PROCESSOR}变量。前者包含当前运行的CPU在CMake的名称，而后者将包含当前正在为其构建的CPU的名称。这是一个细微的差别，在交叉编译时起着非常重要的作用。我们将在第13章，看到更多关于交叉编译的内容。另一种让CMake检测主机处理器体系结构，是使用\texttt{C}或\texttt{C++中}定义的符号，结合CMake的\texttt{try\_run}函数，尝试构建执行的源代码(见第5.8节)分支的预处理符号。这将返回已定义错误码，这些错误可以在CMake端捕获(此策略的灵感来自 \texthref{https://github.com/axr/cmake/blob/master/targetarch.cmake} ):

\begin{codebox}[cpp]
#if defined(__i386) || defined(__i386__) || defined(_M_IX86)
    #error cmake_arch i386
#elif defined(__x86_64) || defined(__x86_64__) || defined(__amd64) || defined(_M_X64)
    #error cmake_arch x86_64
#endif
\end{codebox}

这种策略也是检测目标处理器体系结构的推荐策略，因为CMake似乎没有提供可移植的内在解决方案。另一种选择，将只使用CMake，完全不使用预处理器，代价是为每种情况设置不同的源文件，然后使用\texttt{target\_source }命令将其设置为可执行目标\texttt{arch-dependent}依赖的源文件:

\begin{codebox}[cmake]
add_executable(arch-dependent "")

if(CMAKE_HOST_SYSTEM_PROCESSOR MATCHES "i386")
    message(STATUS "i386 architecture detected")
    target_sources(arch-dependent
        PRIVATE
        arch-dependent-i386.cpp
    )
elseif(CMAKE_HOST_SYSTEM_PROCESSOR MATCHES "i686")
    message(STATUS "i686 architecture detected")
    target_sources(arch-dependent
        PRIVATE
            arch-dependent-i686.cpp
    )
elseif(CMAKE_HOST_SYSTEM_PROCESSOR MATCHES "x86_64")
    message(STATUS "x86_64 architecture detected")
    target_sources(arch-dependent
        PRIVATE
            arch-dependent-x86_64.cpp
    )
else()
    message(STATUS "host processor architecture is unknown")
endif()
\end{codebox}

这种方法，显然需要对现有项目进行更多的工作，因为源文件需要分离。此外，不同源文件之间的代码复制肯定也会成为问题。


% --- recipe 5: 2.5-chinese.md ---
\section{2.5 检测处理器指令集}

\begin{myInfoBox}
\textbf{NOTE}：此示例代码可以在 \texthref{https://github.com/dev-cafe/cmake-cookbook/tree/v1.0/chapter-02/recipe-05} 中找到，包含一个C++示例。该示例在CMake 3.10版(或更高版本)中是有效的，并且已经在GNU/Linux、macOS和Windows上进行过测试。
\end{myInfoBox}

本示例中，我们将讨论如何在CMake的帮助下检测主机处理器支持的指令集。这个功能是较新版本添加到CMake中的，需要CMake 3.10或更高版本。检测到的主机系统信息，可用于设置相应的编译器标志，或实现可选的源代码编译，或根据主机系统生成源代码。本示例中，我们的目标是检测主机系统信息，使用预处理器定义将其传递给\texttt{C++}源代码，并将信息打印到输出中。

\mySmallsection{准备工作}

我们是\texttt{C++}源码(\texttt{processor-info.cpp})如下所示：

\begin{codebox}[cpp]
#include "config.h"

#include <cstdlib>
#include <iostream>

int main()
{
  std::cout << "Number of logical cores: "
            << NUMBER_OF_LOGICAL_CORES << std::endl;
  std::cout << "Number of physical cores: "
            << NUMBER_OF_PHYSICAL_CORES << std::endl;
  std::cout << "Total virtual memory in megabytes: "
            << TOTAL_VIRTUAL_MEMORY << std::endl;
  std::cout << "Available virtual memory in megabytes: "
            << AVAILABLE_VIRTUAL_MEMORY << std::endl;
  std::cout << "Total physical memory in megabytes: "
            << TOTAL_PHYSICAL_MEMORY << std::endl;
  std::cout << "Available physical memory in megabytes: "
            << AVAILABLE_PHYSICAL_MEMORY << std::endl;
  std::cout << "Processor is 64Bit: "
            << IS_64BIT << std::endl;
  std::cout << "Processor has floating point unit: "
            << HAS_FPU << std::endl;
  std::cout << "Processor supports MMX instructions: "
            << HAS_MMX << std::endl;
  std::cout << "Processor supports Ext. MMX instructions: "
            << HAS_MMX_PLUS << std::endl;
  std::cout << "Processor supports SSE instructions: "
            << HAS_SSE << std::endl;
  std::cout << "Processor supports SSE2 instructions: "
            << HAS_SSE2 << std::endl;
  std::cout << "Processor supports SSE FP instructions: "
            << HAS_SSE_FP << std::endl;
  std::cout << "Processor supports SSE MMX instructions: "
            << HAS_SSE_MMX << std::endl;
  std::cout << "Processor supports 3DNow instructions: "
            << HAS_AMD_3DNOW << std::endl;
  std::cout << "Processor supports 3DNow+ instructions: "
            << HAS_AMD_3DNOW_PLUS << std::endl;
  std::cout << "IA64 processor emulating x86 : "
            << HAS_IA64 << std::endl;
  std::cout << "OS name: "
            << OS_NAME << std::endl;
  std::cout << "OS sub-type: "
            << OS_RELEASE << std::endl;
  std::cout << "OS build ID: "
            << OS_VERSION << std::endl;
  std::cout << "OS platform: "
            << OS_PLATFORM << std::endl;
  return EXIT_SUCCESS;
}
\end{codebox}

其包含\texttt{config.h}头文件，我们将使用\texttt{config.h.in}生成这个文件。\texttt{config.h.in}如下:

\begin{codebox}[cpp]
#pragma once

#define NUMBER_OF_LOGICAL_CORES @_NUMBER_OF_LOGICAL_CORES@
#define NUMBER_OF_PHYSICAL_CORES @_NUMBER_OF_PHYSICAL_CORES@
#define TOTAL_VIRTUAL_MEMORY @_TOTAL_VIRTUAL_MEMORY@
#define AVAILABLE_VIRTUAL_MEMORY @_AVAILABLE_VIRTUAL_MEMORY@
#define TOTAL_PHYSICAL_MEMORY @_TOTAL_PHYSICAL_MEMORY@
#define AVAILABLE_PHYSICAL_MEMORY @_AVAILABLE_PHYSICAL_MEMORY@
#define IS_64BIT @_IS_64BIT@
#define HAS_FPU @_HAS_FPU@
#define HAS_MMX @_HAS_MMX@
#define HAS_MMX_PLUS @_HAS_MMX_PLUS@
#define HAS_SSE @_HAS_SSE@
#define HAS_SSE2 @_HAS_SSE2@
#define HAS_SSE_FP @_HAS_SSE_FP@
#define HAS_SSE_MMX @_HAS_SSE_MMX@
#define HAS_AMD_3DNOW @_HAS_AMD_3DNOW@
#define HAS_AMD_3DNOW_PLUS @_HAS_AMD_3DNOW_PLUS@
#define HAS_IA64 @_HAS_IA64@
#define OS_NAME "@_OS_NAME@"
#define OS_RELEASE "@_OS_RELEASE@"
#define OS_VERSION "@_OS_VERSION@"
#define OS_PLATFORM "@_OS_PLATFORM@"
\end{codebox}

\mySmallsection{如何实施}

我们将使用CMake为平台填充\texttt{config.h}中的定义，并将示例源文件编译为可执行文件:

\begin{enumerate}
\item 首先，我们定义了CMake最低版本、项目名称和项目语言:

\begin{codebox}[cmake]
cmake_minimum_required(VERSION 3.10 FATAL_ERROR)
project(recipe-05 CXX)
\end{codebox}

\item 然后，定义目标可执行文件及其源文件，并包括目录:

\begin{codebox}[cmake]
add_executable(processor-info "")
target_sources(processor-info
  PRIVATE
    processor-info.cpp
  )
target_include_directories(processor-info
  PRIVATE
      ${PROJECT_BINARY_DIR}
  )
\end{codebox}

\item 继续查询主机系统的信息，获取一些关键字:

\begin{codebox}[cmake]
foreach(key
  IN ITEMS
    NUMBER_OF_LOGICAL_CORES
    NUMBER_OF_PHYSICAL_CORES
    TOTAL_VIRTUAL_MEMORY
    AVAILABLE_VIRTUAL_MEMORY
    TOTAL_PHYSICAL_MEMORY
    AVAILABLE_PHYSICAL_MEMORY
    IS_64BIT
    HAS_FPU
    HAS_MMX
    HAS_MMX_PLUS
    HAS_SSE
    HAS_SSE2
    HAS_SSE_FP
    HAS_SSE_MMX
    HAS_AMD_3DNOW
    HAS_AMD_3DNOW_PLUS
    HAS_IA64
    OS_NAME
    OS_RELEASE
    OS_VERSION
    OS_PLATFORM
  )
  cmake_host_system_information(RESULT _${key} QUERY ${key})
endforeach()
\end{codebox}

\item 定义了相应的变量后，配置\texttt{config.h}:

\begin{codebox}[cmake]
configure_file(config.h.in config.h @ONLY)
\end{codebox}

\item 现在准备好配置、构建和测试项目:

\begin{codebox}[shell]
$ mkdir -p build
$ cd build
$ cmake ..
$ cmake --build .
$ ./processor-info
Number of logical cores: 4
Number of physical cores: 2
Total virtual memory in megabytes: 15258
Available virtual memory in megabytes: 14678
Total physical memory in megabytes: 7858
Available physical memory in megabytes: 4072
Processor is 64Bit: 1
Processor has floating point unit: 1
Processor supports MMX instructions: 1
Processor supports Ext. MMX instructions: 0
Processor supports SSE instructions: 1
Processor supports SSE2 instructions: 1
Processor supports SSE FP instructions: 0
Processor supports SSE MMX instructions: 0
Processor supports 3DNow instructions: 0
Processor supports 3DNow+ instructions: 0
IA64 processor emulating x86 : 0
OS name: Linux
OS sub-type: 4.16.7-1-ARCH
OS build ID: #1 SMP PREEMPT Wed May 2 21:12:36 UTC 2018
OS platform: x86_64
\end{codebox}

\item 输出会随着处理器的不同而变化。

\end{enumerate}

\mySmallsection{工作原理}

\texttt{CMakeLists.txt}中的\texttt{foreach}循环会查询多个键值，并定义相应的变量。此示例的核心函数是\texttt{cmake\_host\_system\_information}，它查询运行CMake的主机系统的系统信息。本例中，我们对每个键使用了一个函数调用。然后，使用这些变量来配置\texttt{config.h.in}中的占位符，输入并生成\texttt{config.h}。此配置使用\texttt{configure\_file}命令完成。最后，\texttt{config.h}包含在\texttt{processor-info.cpp}中。编译后，它将把值打印到屏幕上。我们将在第5章(配置时和构建时操作)和第6章(生成源代码)中重新讨论这种方法。

\mySmallsection{更多信息}

对于更细粒度的处理器指令集检测，请考虑以下模块: \texthref{https://github.com/VcDevel/Vc/blob/master/cmake/OptimizeForArchitecture.cmake} 。有时候，构建代码的主机可能与运行代码的主机不一样。在计算集群中，登录节点的体系结构可能与计算节点上的体系结构不同。解决此问题的一种方法是，将配置和编译作为计算步骤，提交并部署到相应计算节点上。


% --- recipe 6: 2.6-chinese.md ---
\section{2.6 为Eigen库使能向量化}

\begin{myInfoBox}
\textbf{NOTE}：此示例代码可以在 \texthref{https://github.com/dev-cafe/cmake-cookbook/tree/v1.0/chapter-02/recipe-06} 中找到，包含一个C++示例。该示例在CMake 3.5版(或更高版本)中是有效的，并且已经在GNU/Linux、macOS和Windows上进行过测试。
\end{myInfoBox}

处理器的向量功能，可以提高代码的性能。对于某些类型的运算来说尤为甚之，例如：线性代数。本示例将展示如何使能矢量化，以便使用线性代数的Eigen C++库加速可执行文件。

\mySmallsection{准备工作}

我们用Eigen C++模板库，用来进行线性代数计算，并展示如何设置编译器标志来启用向量化。这个示例的源代码\texttt{linear-algebra.cpp}文件:

\begin{codebox}[cpp]
#include <chrono>
#include <iostream>

#include <Eigen/Dense>

EIGEN_DONT_INLINE
double simple_function(Eigen::VectorXd &va, Eigen::VectorXd &vb)
{
  // this simple function computes the dot product of two vectors
  // of course it could be expressed more compactly
  double d = va.dot(vb);
  return d;
}

int main()
{
  int len = 1000000;
  int num_repetitions = 100;
  
  // generate two random vectors
  Eigen::VectorXd va = Eigen::VectorXd::Random(len);
  Eigen::VectorXd vb = Eigen::VectorXd::Random(len);
  
  double result;
  auto start = std::chrono::system_clock::now();
  for (auto i = 0; i < num_repetitions; i++)
  {
    result = simple_function(va, vb);
  }
  auto end = std::chrono::system_clock::now();
  auto elapsed_seconds = end - start;
  
  std::cout << "result: " << result << std::endl;
  std::cout << "elapsed seconds: " << elapsed_seconds.count() << std::endl;
}
\end{codebox}

我们期望向量化可以加快\texttt{simple\_function}中的点积操作。

\mySmallsection{如何实施}

根据Eigen库的文档，设置适当的编译器标志就足以生成向量化的代码。让我们看看\texttt{CMakeLists.txt}:

\begin{enumerate}
\item 声明一个\texttt{C++11}项目:

\begin{codebox}[cmake]
cmake_minimum_required(VERSION 3.5 FATAL_ERROR)
project(recipe-06 LANGUAGES CXX)
set(CMAKE_CXX_STANDARD 11)
set(CMAKE_CXX_EXTENSIONS OFF)
set(CMAKE_CXX_STANDARD_REQUIRED ON)
\end{codebox}

\item 使用Eigen库，我们需要在系统上找到它的头文件:

\begin{codebox}[cmake]
find_package(Eigen3 3.3 REQUIRED CONFIG)
\end{codebox}

\item \texttt{CheckCXXCompilerFlag.cmake}标准模块文件:

\begin{codebox}[cmake]
include(CheckCXXCompilerFlag)
\end{codebox}

\item 检查\texttt{ -march=native }编译器标志是否工作:

\begin{codebox}[cmake]
check_cxx_compiler_flag("-march=native" _march_native_works)
\end{codebox}

\item 另一个选项\texttt{-xHost}编译器标志也开启:

\begin{codebox}[cmake]
check_cxx_compiler_flag("-xHost" _xhost_works)
\end{codebox}

\item 设置了一个空变量\texttt{\_CXX\_FLAGS}，来保存刚才检查的两个编译器中找到的编译器标志。如果看到\texttt{\_march\_native\_works}，我们将\texttt{\_CXX\_FLAGS}设置为\texttt{-march=native}。如果看到\texttt{\_xhost\_works}，我们将\texttt{\_CXX\_FLAGS}设置为\texttt{-xHost}。如果它们都不起作用，\texttt{\_CXX\_FLAGS}将为空，并禁用矢量化:

\begin{codebox}[cmake]
set(_CXX_FLAGS)
if(_march_native_works)
    message(STATUS "Using processor's vector instructions (-march=native compiler flag set)")
    set(_CXX_FLAGS "-march=native")
elseif(_xhost_works)
    message(STATUS "Using processor's vector instructions (-xHost compiler flag set)")
    set(_CXX_FLAGS "-xHost")
else()
    message(STATUS "No suitable compiler flag found for vectorization")
endif()
\end{codebox}

\item 为了便于比较，我们还为未优化的版本定义了一个可执行目标，不使用优化标志:

\begin{codebox}[cmake]
add_executable(linear-algebra-unoptimized linear-algebra.cpp)
target_link_libraries(linear-algebra-unoptimized
  PRIVATE
    Eigen3::Eigen
  )
\end{codebox}

\item 此外，我们定义了一个优化版本:

\begin{codebox}[cmake]
add_executable(linear-algebra linear-algebra.cpp)
target_compile_options(linear-algebra
  PRIVATE
    ${_CXX_FLAGS}
  )
target_link_libraries(linear-algebra
  PRIVATE
    Eigen3::Eigen
  )
\end{codebox}

\item 让我们比较一下这两个可执行文件——首先我们配置(在本例中，\texttt{-march=native\_works}):

\begin{codebox}[shell]
$ mkdir -p build
$ cd build
$ cmake ..
...
-- Performing Test _march_native_works
-- Performing Test _march_native_works - Success
-- Performing Test _xhost_works
-- Performing Test _xhost_works - Failed
-- Using processor's vector instructions (-march=native compiler flag set)
...
\end{codebox}

\item 最后，让我们编译可执行文件，并比较运行时间:

\begin{codebox}[shell]
$ cmake --build .
$ ./linear-algebra-unoptimized
result: -261.505
elapsed seconds: 1.97964
$ ./linear-algebra
result: -261.505
elapsed seconds: 1.05048
\end{codebox}

\end{enumerate}

\mySmallsection{工作原理}

大多数处理器提供向量指令集，代码可以利用这些特性，获得更高的性能。由于线性代数运算可以从Eigen库中获得很好的加速，所以在使用Eigen库时，就要考虑向量化。我们所要做的就是，指示编译器为我们检查处理器，并为当前体系结构生成本机指令。不同的编译器供应商会使用不同的标志来实现这一点：GNU编译器使用\texttt{-march=native}标志来实现这一点，而Intel编译器使用\texttt{-xHost}标志。使用\texttt{ CheckCXXCompilerFlag.cmake}模块提供的\texttt{check\_cxx\_compiler\_flag}函数进行编译器标志的检查:

\texttt{check\_cxx\_compiler\_flag("-march=native" \_march\_native\_works)}

这个函数接受两个参数:

\begin{itemize}
\item 第一个是要检查的编译器标志。

\item 第二个是用来存储检查结果(true或false)的变量。如果检查为真，我们将工作标志添加到\texttt{\_CXX\_FLAGS}变量中，该变量将用于为可执行目标设置编译器标志。

\end{itemize}

\mySmallsection{更多信息}

本示例可与前一示例相结合，可以使用\texttt{cmake\_host\_system\_information}查询处理器功能。


